\documentclass[11pt]{article}

\usepackage[utf8]{inputenc}
\usepackage[T1]{fontenc}
\usepackage{lmodern}
\usepackage{geometry}
\usepackage{amsmath,amssymb,amsfonts,amsthm,mathtools}
\usepackage{bm}
\usepackage{enumitem}
\usepackage{microtype}
\usepackage{hyperref}
\usepackage{titlesec}
\usepackage{booktabs}
\usepackage{array}
\usepackage{setspace}
\usepackage{mathrsfs}
\usepackage{physics}

\geometry{margin=1in}
\hypersetup{colorlinks=true, linkcolor=black, urlcolor=blue, citecolor=black}
\setlist[enumerate]{topsep=0.4em,itemsep=0.35em}
\setlist[itemize]{topsep=0.3em,itemsep=0.25em}
\titleformat{\section}{\large\bfseries}{\thesection.}{0.5em}{}
\titleformat{\subsection}{\normalsize\bfseries}{\thesubsection.}{0.5em}{}
\setstretch{1.06}

\newcommand{\Aether}{\AE{}ther}
\newcommand{\Aflow}{\AE{}ther-flow}
\newcommand{\TheoryName}{\Aflow{} Interpretation of Relativity}
\newcommand{\TheoryProgram}{\Aether{} / \Aflow{} framework}

\newcommand{\CompletedMetric}{g^{\sharp}}
\newcommand{\MatterSpace}{\Psi}

\newcommand{\Car}{\mathrm{car}}
\newcommand{\Cur}{\mathrm{cur}}
\newcommand{\Hom}{\mathrm{hom}}

\newcommand{\ForSpace}[1]{\mathcal I_{#1}^{\mathrm{for}}}
\newcommand{\PrimShell}{\mathcal S_{\mathrm{prim}}}

\newcommand{\Reduction}[1]{\mathcal R_{#1}}
\newcommand{\CurrentCarrier}[1]{\mathfrak C_{#1}^{\mathrm{cur}}}
\newcommand{\EqCarrier}[1]{\mathfrak C_{#1}^{\mathrm{eq}}}

\newcommand{\MetricOut}{\mathscr G_{\mathrm{out}}}
\newcommand{\MatterOut}{\mathscr T_{\mathrm{out}}}

\newcommand{\DataSpace}[1]{\mathcal Y_{#1}^{\mathrm{obs}}}
\newcommand{\AmbientTarget}[1]{E_{#1}^{\mathrm{amb}}}

\newcommand{\SubSpace}[1]{\mathcal M_{#1}^{\mathrm{sub}}}
\newcommand{\SubBody}[1]{\mathcal B_{#1}^{\mathrm{sub}}}
\newcommand{\SubTarget}[1]{E_{#1}^{\mathrm{sub}}}
\newcommand{\SubPair}[1]{\mathfrak s_{#1}^{\mathrm{sub}}}

\newcommand{\LiftSpace}[1]{\mathcal L_{#1}^{\mathrm{lift}}}
\newcommand{\LiftBody}[1]{\mathcal B_{#1}^{\mathrm{lift}}}
\newcommand{\PrecSpace}[1]{\mathcal P_{#1}^{\mathrm{prec}}}
\newcommand{\PrecBody}[1]{\mathcal B_{#1}^{\mathrm{prec}}}
\newcommand{\OrigSpace}[1]{\mathcal O_{#1}^{\mathrm{orig}}}
\newcommand{\OrigBody}[1]{\mathcal B_{#1}^{\mathrm{orig}}}

\newcommand{\GroundSpace}[1]{\mathcal G_{#1}^{\mathrm{grd}}}
\newcommand{\GroundBody}[1]{\mathcal B_{#1}^{\mathrm{grd}}}
\newcommand{\GroundProj}[1]{\Pi_{#1}^{\mathrm{grd}}}

\newcommand{\FoundSpace}[1]{\mathcal F_{#1}^{\mathrm{fdn}}}
\newcommand{\FoundBody}[1]{\mathcal B_{#1}^{\mathrm{fdn}}}
\newcommand{\FoundProj}[1]{\Pi_{#1}^{\mathrm{fdn}}}
\newcommand{\FoundGroundMin}[1]{\mathfrak f_{#1}^{\mathrm{fdn,grd}}}
\newcommand{\FoundOrigMin}[1]{\mathfrak f_{#1}^{\mathrm{fdn,orig}}}
\newcommand{\FoundPrecMin}[1]{\mathfrak f_{#1}^{\mathrm{fdn,prec}}}
\newcommand{\FoundLiftMin}[1]{\mathfrak f_{#1}^{\mathrm{fdn,lift}}}
\newcommand{\FoundSubMin}[1]{\mathfrak m_{#1}^{\mathrm{fdn}}}
\newcommand{\FoundSection}[1]{\Gamma_{#1}^{\mathrm{fdn}}}
\newcommand{\FoundFunctional}[1]{\mathscr W_{#1}^{\mathrm{fdn}}}
\newcommand{\FoundData}[1]{\Lambda_{#1}^{\mathrm{fdn,dat}}}
\newcommand{\FoundShellMap}[1]{\Lambda_{#1}^{\mathrm{fdn,sh}}}
\newcommand{\FoundShellSet}[1]{\mathcal Z_{#1}^{\mathrm{fdn}}}
\newcommand{\FoundSym}[1]{\mathbb G_{#1}^{\mathrm{fdn}}}
\newcommand{\FoundTime}[1]{\vartheta_{#1}^{\mathrm{fdn}}}
\newcommand{\FoundChart}[1]{\Xi_{#1}^{\mathrm{fdn}}}
\newcommand{\FoundSupport}[1]{\Theta_{#1}^{\mathrm{fdn,sup}}}
\newcommand{\FoundSupportShell}[1]{\mathcal Z_{#1}^{\mathrm{fdn,sup}}}
\newcommand{\FoundZeroCore}[1]{\mathcal Z_{#1}^{0,\mathrm{fdn}}}
\newcommand{\FoundEqChart}[1]{\Psi_{#1}^{\mathrm{fdn,sup}}}
\newcommand{\FoundZeroChart}[1]{\Psi_{#1}^{0,\mathrm{fdn}}}
\newcommand{\FoundSplit}[1]{\Phi_{#1}^{\mathrm{fdn}}}
\newcommand{\FoundCharge}[1]{\mathcal Q_{#1}^{\mathrm{fdn}}}
\newcommand{\FoundResidual}[1]{\mathcal R_{#1}^{\mathrm{fdn}}}
\newcommand{\FoundEmbed}[1]{\zeta_{#1}^{\mathrm{fdn}}}
\newcommand{\FoundKernelCh}[1]{\widetilde K_{#1}^{0,\mathrm{fdn}}}
\newcommand{\FoundStateComp}[1]{P_{#1}^{\mathrm{fdn,co}}}

\newcommand{\RootSpace}[1]{\mathcal R_{#1}^{\mathrm{rt}}}
\newcommand{\RootBody}[1]{\mathcal B_{#1}^{\mathrm{rt}}}
\newcommand{\RootProj}[1]{\Pi_{#1}^{\mathrm{rt}}}
\newcommand{\RootFoundMin}[1]{\mathfrak f_{#1}^{\mathrm{rt,fdn}}}
\newcommand{\RootGroundMin}[1]{\mathfrak f_{#1}^{\mathrm{rt,grd}}}
\newcommand{\RootOrigMin}[1]{\mathfrak f_{#1}^{\mathrm{rt,orig}}}
\newcommand{\RootPrecMin}[1]{\mathfrak f_{#1}^{\mathrm{rt,prec}}}
\newcommand{\RootLiftMin}[1]{\mathfrak f_{#1}^{\mathrm{rt,lift}}}
\newcommand{\RootSubMin}[1]{\mathfrak m_{#1}^{\mathrm{rt}}}
\newcommand{\RootSection}[1]{\Gamma_{#1}^{\mathrm{rt}}}
\newcommand{\RootFunctional}[1]{\mathscr W_{#1}^{\mathrm{rt}}}
\newcommand{\RootData}[1]{\Lambda_{#1}^{\mathrm{rt,dat}}}
\newcommand{\RootShellMap}[1]{\Lambda_{#1}^{\mathrm{rt,sh}}}
\newcommand{\RootSym}[1]{\mathbb G_{#1}^{\mathrm{rt}}}
\newcommand{\RootTime}[1]{\vartheta_{#1}^{\mathrm{rt}}}
\newcommand{\RootShellSet}[1]{\mathcal Z_{#1}^{\mathrm{rt}}}
\newcommand{\RootChart}[1]{\Xi_{#1}^{\mathrm{rt}}}
\newcommand{\RootSupport}[1]{\Theta_{#1}^{\mathrm{rt,sup}}}
\newcommand{\RootSupportShell}[1]{\mathcal Z_{#1}^{\mathrm{rt,sup}}}
\newcommand{\RootZeroCore}[1]{\mathcal Z_{#1}^{0,\mathrm{rt}}}
\newcommand{\RootEqChart}[1]{\Psi_{#1}^{\mathrm{rt,sup}}}
\newcommand{\RootZeroChart}[1]{\Psi_{#1}^{0,\mathrm{rt}}}
\newcommand{\RootSplit}[1]{\Phi_{#1}^{\mathrm{rt}}}
\newcommand{\RootCharge}[1]{\mathcal Q_{#1}^{\mathrm{rt}}}
\newcommand{\RootResidual}[1]{\mathcal R_{#1}^{\mathrm{rt}}}
\newcommand{\RootEmbed}[1]{\zeta_{#1}^{\mathrm{rt}}}
\newcommand{\RootKernelCh}[1]{\widetilde K_{#1}^{0,\mathrm{rt}}}
\newcommand{\RootStateComp}[1]{P_{#1}^{\mathrm{rt,co}}}

\newcommand{\ResSpace}[1]{\mathcal U_{#1}^{\mathrm{sup,res}}}
\newcommand{\ResTarget}[1]{S_{#1}^{\mathrm{sup,res}}}
\newcommand{\SeedHidden}[1]{\mathcal H_{#1}^{\mathrm{sup}}}
\newcommand{\SeedHiddenBody}[1]{D_{#1}^{\mathrm{sup}}}
\newcommand{\HiddenChart}[1]{\Psi_{#1}^{\mathrm{hid}}}

\newcommand{\EqnSpace}[1]{\mathcal U_{#1}^{\mathrm{eqn}}}
\newcommand{\EqnTarget}[1]{Q_{#1}^{\mathrm{eqn}}}
\newcommand{\EqnPair}[1]{\mathfrak b_{#1}^{\mathrm{eqn}}}
\newcommand{\EqnData}[1]{\Lambda_{#1}^{\mathrm{eqn,dat}}}
\newcommand{\EqnShell}[1]{\Lambda_{#1}^{\mathrm{eqn,sh}}}
\newcommand{\EqnHidden}[1]{H_{#1}^{\mathrm{eqn}}}
\newcommand{\EqnZero}[1]{V_{#1}^{\mathrm{eqn},0}}

\newcommand{\RootCore}[1]{\mathfrak F_{#1}^{\mathrm{rt,core}}}
\newcommand{\RootMapFound}[1]{\mathcal E_{#1}^{\mathrm{fdn}}}
\newcommand{\RootMapGround}[1]{\mathcal E_{#1}^{\mathrm{grd}}}
\newcommand{\RootMapOrig}[1]{\mathcal E_{#1}^{\mathrm{orig}}}
\newcommand{\RootMapPrec}[1]{\mathcal E_{#1}^{\mathrm{prec}}}
\newcommand{\RootMapLift}[1]{\mathcal E_{#1}^{\mathrm{lift}}}
\newcommand{\RootMapSub}[1]{\mathcal E_{#1}^{\mathrm{sub}}}
\newcommand{\RootMapShell}[1]{\mathcal E_{#1}^{\mathrm{sh}}}
\newcommand{\RootMapSupport}[1]{\mathcal E_{#1}^{\mathrm{sup}}}
\newcommand{\RootMapZero}[1]{\mathcal E_{#1}^{0}}

\newtheorem{definition}{Definition}
\newtheorem{proposition}{Proposition}
\newtheorem{remark}{Remark}
\newtheorem{corollary}{Corollary}
\newtheorem{theorem}{Theorem}

\title{\textbf{The \TheoryName{}}\\[0.4em]
\large Primitive-Reservoir Observer-Reduction\\
\large Foundation-Generating Substrate Root Dynamics\\
\large Necessity / Uniqueness Theorem}
\author{}
\date{}

\begin{document}

\maketitle

\begin{abstract}
The recorded theorem deriving ground-generating substrate foundation dynamics from foundation-generating substrate root dynamics already pushed the active line one layer below the settled foundation frontier. Under the routing safeguard recorded at that frontier, however, that derivation theorem remained only a same-output deeper relay: it preserved exactly the same foundation package and therefore did not by itself supersede the benchmark-facing status of the foundation-core compression theorem. If the root layer is to become a genuine benchmark-facing gain, it must now do more than reproduce the same lower package. It must remove benchmark-relevant freedom inside the root layer itself.

The present note proves exactly that sharper result. For each sector it isolates the benchmark-relevant \emph{fiber-minimizing exact-shell root core} of a root-faithful presentation: the root-to-foundation projection together with the six recorded root fiber-minimizer images, the exact root shell and its observer-data and shell-response readouts, the shell chart to the same reservoir space, the support shell and zero core, the zero/hidden split in the same equation variables, the hidden charge and exact root zero-response realization, and the fixed-data solvability and coercive control carried on that core. The theorem then proves that every root-faithful presentation inducing the fixed foundation package determines exactly the same root core up to canonical root-faithful equivalence.

The benchmark boundary remains unchanged throughout. The one operative metric is still exactly \(\CompletedMetric\), the ordinary matter route is unchanged, there is no second operative metric, and the low-energy matter law is not enlarged. The positive result is therefore disciplined and benchmark facing: the still-more-primitive root layer is no longer merely available, but forced up to root-faithful equivalence. The honest next burden is deeper again: derive or justify the root dynamics from still more primitive substrate structure, or prove a genuinely smaller theorem-level collapse strictly inside that root core.
\end{abstract}

\section{Introduction}

The active derivational program of \TheoryName{} now has a cleaner task at root scope than another same-output relay. The foundation necessity / uniqueness theorem already spent the honest internal compression move available at the previous frontier, and the later root derivation theorem then showed that the fixed foundation package can indeed be recovered from a deeper root layer. That deeper theorem matters because it removes the foundation package as the deepest recorded stopping point. But by the current routing safeguard it does not yet count as the new benchmark-facing frontier on its own.

Why not? Because depth alone is no longer enough. Once a manuscript preserves the same benchmark one-metric package while merely relocating the unexplained substrate layer deeper, it changes project status only if it also supplies a necessity, uniqueness, compression, or comparably sharp result at that deeper level. The honest next move below the foundation frontier is therefore not another root-side restatement of shell, support, reservoir, or hidden-sector data. It is to ask whether the benchmark-relevant part of the root package itself is already forced once the fixed foundation package is required.

That is the point of the present note. It does not claim a completed first-principles derivation of gravity from final substrate microphysics. Its gain is narrower and more honest. It shows that, at the present root level, the part of the root package that actually matters for the active one-metric derivational line is no longer a free choice. Once one requires a root presentation to induce exactly the fixed foundation package, the benchmark-relevant root core is forced up to canonical equivalence.

\paragraph{Framework and claim boundary.}
\TheoryProgram{} denotes the broader \Aether{} / \Aflow{} framework built on the claims that the \Aether{} is the underlying four-dimensional substrate of reality, the \Aflow{} is its intrinsic ordered motion, observed three-dimensional space is the local experiential slice of that deeper substrate, S-time is the experienced order of change arising from matter, light, and the \Aflow{}, and observed expansion is the three-dimensional appearance of deeper four-dimensional ordered motion. Within that framework, \TheoryName{} denotes the exact-closure theory in which the effective gravitational dynamics are adopted to be exactly Einsteinian with universal matter coupling. In the active sequence, \emph{adoption} means use of that established relativistic dynamics without claiming substrate derivation, while \emph{derivation} is reserved for a first-principles recovery from explicit substrate variables.

The present note stays strictly inside that boundary. It does not alter the benchmark exact-closure package. It does not introduce a second operative metric or a new low-energy matter law. It only proves that the current root layer is already rigid at benchmark scope once it is required to reproduce the fixed foundation package.

\section{Fixed Inputs}

\begin{definition}[Recorded primitive continuation data]
\label{def:fixed}
Fix the following continuation data from the active primitive-reservoir observer-reduction line.
\begin{enumerate}
    \item For each sector label \(\sigma\in\{\Car,\Cur,\Hom\}\), a primitive forcing shell
    \[
    \ForSpace{\sigma},
    \]
    a visible reduction map
    \[
    \Reduction{\sigma}:\ForSpace{\sigma}\to X_{\sigma},
    \]
    and shell-level current and equilibrium carriers
    \[
    \CurrentCarrier{\sigma}:\ForSpace{\sigma}\to Z_{\sigma}^{\mathrm{cur}},
    \qquad
    \EqCarrier{\sigma}:\ForSpace{\sigma}\to Z_{\sigma}^{\mathrm{eq}}.
    \]
    \item The product shell
    \[
    \PrimShell
    \equiv
    \ForSpace{\Car}\times \ForSpace{\Cur}\times \ForSpace{\Hom}.
    \]
    \item The recorded metric and ordinary-matter outputs
    \[
    \MetricOut:\PrimShell\to\{\text{observer metrics}\},
    \qquad
    \MatterOut:\PrimShell\times\MatterSpace\to\{\text{observer stress tensors}\},
    \]
    satisfying
    \begin{equation}
    \MetricOut(\mathbf w)=\CompletedMetric,
    \qquad
    \MatterOut(\mathbf w,\psi)
    =
    T_{\mu\nu}^{\mathrm{matter}}[\CompletedMetric,\psi]
    \label{eq:fixedoutputs}
    \end{equation}
    for every \(\mathbf w\in\PrimShell\) and every matter configuration \(\psi\in\MatterSpace\).
\end{enumerate}
\end{definition}

\begin{definition}[Recorded foundation target]
\label{def:foundation}
Fix \(\sigma\in\{\Car,\Cur,\Hom\}\). The fixed ground-generating substrate foundation package on the active line consists of the following data.
\begin{enumerate}
    \item A finite-dimensional Euclidean foundation state space
    \[
    \FoundSpace{\sigma},
    \]
    a compact convex foundation body
    \[
    \FoundBody{\sigma}\subset\FoundSpace{\sigma}
    \]
    with nonempty interior, an affine surjection
    \[
    \FoundProj{\sigma}:\FoundSpace{\sigma}\to\GroundSpace{\sigma},
    \]
    and smooth foundation fiber minimizers
    \[
    \FoundGroundMin{\sigma},\quad
    \FoundOrigMin{\sigma},\quad
    \FoundPrecMin{\sigma},\quad
    \FoundLiftMin{\sigma},\quad
    \FoundSubMin{\sigma}
    \]
    on the fixed ground, origin, precursor, lift, and substrate fibers of the active line.
    \item The same substrate target \(\SubTarget{\sigma}\), the same pairing \(\SubPair{\sigma}\), a smooth foundation section
    \[
    \FoundSection{\sigma}:\FoundSpace{\sigma}\to\SubTarget{\sigma},
    \]
    the foundation functional
    \[
    \FoundFunctional{\sigma}(f)
    \equiv
    \SubPair{\sigma}\bigl(\FoundSection{\sigma}(f),\FoundSection{\sigma}(f)\bigr),
    \]
    and the exact foundation shell
    \[
    \FoundShellSet{\sigma}
    \equiv
    \bigl(\FoundSection{\sigma}\bigr)^{-1}(0)\cap\FoundBody{\sigma}.
    \]
    \item Affine foundation observer-data and shell-response readouts
    \[
    \FoundData{\sigma}:\FoundSpace{\sigma}\to\DataSpace{\sigma},
    \qquad
    \FoundShellMap{\sigma}:\FoundSpace{\sigma}\to\AmbientTarget{\sigma},
    \]
    together with a compact symmetry group \(\FoundSym{\sigma}\) and an involution \(\FoundTime{\sigma}\) preserving all displayed foundation data.
    \item A Euclidean foundation shell chart
    \[
    \FoundChart{\sigma}:\FoundShellSet{\sigma}\to\ResSpace{\sigma},
    \]
    a shell-internal support recorder
    \[
    \FoundSupport{\sigma}:\FoundShellSet{\sigma}\to\ResTarget{\sigma},
    \]
    the exact foundation support shell
    \[
    \FoundSupportShell{\sigma}
    \equiv
    \bigl(\FoundSupport{\sigma}\bigr)^{-1}(0)\cap\FoundShellSet{\sigma},
    \]
    a closed embedded foundation zero core
    \[
    \FoundZeroCore{\sigma}\subset\FoundSupportShell{\sigma},
    \]
    a support-shell chart
    \[
    \FoundEqChart{\sigma}:\FoundSupportShell{\sigma}\to\EqnSpace{\sigma},
    \]
    a zero chart
    \[
    \FoundZeroChart{\sigma}:\FoundZeroCore{\sigma}\to\EqnZero{\sigma},
    \]
    and a diffeomorphism
    \[
    \FoundSplit{\sigma}:
    \FoundZeroCore{\sigma}\times\SeedHiddenBody{\sigma}
    \xrightarrow{\cong}
    \FoundSupportShell{\sigma}
    \]
    obeying the same zero/hidden split in the same equation variables.
    \item A hidden charge
    \[
    \FoundCharge{\sigma}:\SeedHidden{\sigma}\to\EqnTarget{\sigma}
    \]
    with positive hidden gap, the hidden recorder
    \[
    \FoundResidual{\sigma}\bigl(\FoundSplit{\sigma}(z,\eta)\bigr)
    \equiv
    \FoundCharge{\sigma}(\eta),
    \]
    and a smooth observer-compatible exact foundation zero-response realization
    \[
    \FoundEmbed{\sigma}:\ForSpace{\sigma}\hookrightarrow\FoundZeroCore{\sigma}
    \]
    recovering the fixed visible, current, and equilibrium shell data.
    \item Fixed-data solvability on \(\FoundZeroCore{\sigma}\) and coercive control on foundation data-invisible directions, recorded through the charted kernel \(\FoundKernelCh{\sigma}\) and orthogonal projector \(\FoundStateComp{\sigma}\).
\end{enumerate}
\end{definition}

\begin{remark}
Definitions \ref{def:fixed} and \ref{def:foundation} are not reopened here. The primitive continuation data, the one operative metric, the ordinary matter route, and the recorded foundation package all remain fixed inputs. The present note only addresses whether any benchmark-relevant freedom still remains inside the deeper root layer that induces that foundation package.
\end{remark}

\section{Root-Faithful Presentations}

\begin{definition}[Root-faithful presentation]
\label{def:faithful}
Fix \(\sigma\in\{\Car,\Cur,\Hom\}\). A \emph{root-faithful presentation} of the fixed foundation package consists of:
\begin{enumerate}
    \item a finite-dimensional Euclidean root state space \(\RootSpace{\sigma}\), a compact convex root body \(\RootBody{\sigma}\subset\RootSpace{\sigma}\) with nonempty interior, an affine surjection
    \[
    \RootProj{\sigma}:\RootSpace{\sigma}\to\FoundSpace{\sigma},
    \]
    and smooth root fiber minimizers
    \[
    \RootFoundMin{\sigma},\quad
    \RootGroundMin{\sigma},\quad
    \RootOrigMin{\sigma},\quad
    \RootPrecMin{\sigma},\quad
    \RootLiftMin{\sigma},\quad
    \RootSubMin{\sigma}
    \]
    on the fixed foundation, ground, origin, precursor, lift, and substrate fibers of the active line;
    \item the same substrate target \(\SubTarget{\sigma}\), the same pairing \(\SubPair{\sigma}\), a smooth root section
    \[
    \RootSection{\sigma}:\RootSpace{\sigma}\to\SubTarget{\sigma},
    \]
    the root functional
    \[
    \RootFunctional{\sigma}(r)
    \equiv
    \SubPair{\sigma}\bigl(\RootSection{\sigma}(r),\RootSection{\sigma}(r)\bigr),
    \]
    and the exact root shell
    \[
    \RootShellSet{\sigma}
    \equiv
    \bigl(\RootSection{\sigma}\bigr)^{-1}(0)\cap\RootBody{\sigma};
    \]
    \item affine root observer-data and shell-response readouts
    \[
    \RootData{\sigma}:\RootSpace{\sigma}\to\DataSpace{\sigma},
    \qquad
    \RootShellMap{\sigma}:\RootSpace{\sigma}\to\AmbientTarget{\sigma},
    \]
    both constant on fibers of \(\RootProj{\sigma}\), together with a compact symmetry group \(\RootSym{\sigma}\) and an involution \(\RootTime{\sigma}\) preserving all displayed root data;
    \item a Euclidean root shell chart
    \[
    \RootChart{\sigma}:\RootShellSet{\sigma}\to\ResSpace{\sigma},
    \]
    a shell-internal support recorder
    \[
    \RootSupport{\sigma}:\RootShellSet{\sigma}\to\ResTarget{\sigma},
    \]
    the exact root support shell
    \[
    \RootSupportShell{\sigma}
    \equiv
    \bigl(\RootSupport{\sigma}\bigr)^{-1}(0)\cap\RootShellSet{\sigma},
    \]
    a closed embedded root zero core
    \[
    \RootZeroCore{\sigma}\subset\RootSupportShell{\sigma},
    \]
    a support-shell chart
    \[
    \RootEqChart{\sigma}:\RootSupportShell{\sigma}\to\EqnSpace{\sigma},
    \]
    a zero chart
    \[
    \RootZeroChart{\sigma}:\RootZeroCore{\sigma}\to\EqnZero{\sigma},
    \]
    and a diffeomorphism
    \[
    \RootSplit{\sigma}:
    \RootZeroCore{\sigma}\times\SeedHiddenBody{\sigma}
    \xrightarrow{\cong}
    \RootSupportShell{\sigma}
    \]
    obeying the same zero/hidden split in the same equation variables, with
    \[
    \RootProj{\sigma}\big|_{\RootShellSet{\sigma}},
    \qquad
    \RootProj{\sigma}\big|_{\RootSupportShell{\sigma}},
    \qquad
    \RootProj{\sigma}\big|_{\RootZeroCore{\sigma}}
    \]
    diffeomorphisms onto the fixed loci \(\FoundShellSet{\sigma}\), \(\FoundSupportShell{\sigma}\), and \(\FoundZeroCore{\sigma}\);
    \item a hidden charge
    \[
    \RootCharge{\sigma}:\SeedHidden{\sigma}\to\EqnTarget{\sigma}
    \]
    with positive hidden gap, the hidden recorder
    \[
    \RootResidual{\sigma}\bigl(\RootSplit{\sigma}(z,\eta)\bigr)
    \equiv
    \RootCharge{\sigma}(\eta),
    \]
    and a smooth observer-compatible exact root zero-response realization
    \[
    \RootEmbed{\sigma}:\ForSpace{\sigma}\hookrightarrow\RootZeroCore{\sigma}
    \]
    recovering the fixed visible, current, and equilibrium shell data;
    \item fixed-data solvability on \(\RootZeroCore{\sigma}\) and coercive control on root data-invisible directions, recorded through the charted kernel \(\RootKernelCh{\sigma}\) and orthogonal projector \(\RootStateComp{\sigma}\).
\end{enumerate}
The presentation is \emph{root faithful} if its projected exact-shell transport recovers the fixed foundation package of Definition \ref{def:foundation} exactly, sector by sector, while preserving the benchmark outputs \eqref{eq:fixedoutputs}. Equivalently, the root presentation is required to induce precisely the recorded foundation body, the same foundation fiber-minimizer hierarchy, the same foundation shell and support transport, the same zero/hidden split in the same equation variables, the same hidden charge data, the same exact zero-response realization, and the same solvability/coercive statements, with no second operative metric and no enlarged low-energy matter law.
\end{definition}

\begin{definition}[Fiber-minimizing exact-shell root core]
\label{def:core}
For a root-faithful presentation, the \emph{fiber-minimizing exact-shell root core} \(\RootCore{\sigma}\) is the tuple consisting of:
\begin{enumerate}
    \item the projection \(\RootProj{\sigma}\) together with the images of the six fiber minimizers \(\RootFoundMin{\sigma}\), \(\RootGroundMin{\sigma}\), \(\RootOrigMin{\sigma}\), \(\RootPrecMin{\sigma}\), \(\RootLiftMin{\sigma}\), and \(\RootSubMin{\sigma}\);
    \item the exact root shell \(\RootShellSet{\sigma}\), the root readouts \(\RootData{\sigma}\) and \(\RootShellMap{\sigma}\), the shell chart \(\RootChart{\sigma}\), and the support recorder \(\RootSupport{\sigma}\);
    \item the support shell \(\RootSupportShell{\sigma}\), the zero core \(\RootZeroCore{\sigma}\), the support-shell chart \(\RootEqChart{\sigma}\), the zero chart \(\RootZeroChart{\sigma}\), and the split diffeomorphism \(\RootSplit{\sigma}\);
    \item the hidden charge \(\RootCharge{\sigma}\), the hidden recorder \(\RootResidual{\sigma}\), the exact zero-response realization \(\RootEmbed{\sigma}\), and the solvability/coercive statements carried by \(\RootKernelCh{\sigma}\) and \(\RootStateComp{\sigma}\).
\end{enumerate}
\end{definition}

\begin{definition}[Root-faithful equivalence]
\label{def:equiv}
Let \({}^{(1)}\!\RootCore{\sigma}\) and \({}^{(2)}\!\RootCore{\sigma}\) be the root cores of two root-faithful presentations of the same fixed foundation package. Their canonical comparison maps are defined by the lower data that both cores induce:
\begin{align}
\RootMapFound{\sigma}\bigl({}^{(1)}\!\RootFoundMin{\sigma}(f)\bigr)
&\equiv
{}^{(2)}\!\RootFoundMin{\sigma}(f),
\label{eq:eqfdn}
\\
\RootMapGround{\sigma}\bigl({}^{(1)}\!\RootGroundMin{\sigma}(g)\bigr)
&\equiv
{}^{(2)}\!\RootGroundMin{\sigma}(g),
\label{eq:eqgrd}
\\
\RootMapOrig{\sigma}\bigl({}^{(1)}\!\RootOrigMin{\sigma}(o)\bigr)
&\equiv
{}^{(2)}\!\RootOrigMin{\sigma}(o),
\label{eq:eqorig}
\\
\RootMapPrec{\sigma}\bigl({}^{(1)}\!\RootPrecMin{\sigma}(p)\bigr)
&\equiv
{}^{(2)}\!\RootPrecMin{\sigma}(p),
\label{eq:eqprec}
\\
\RootMapLift{\sigma}\bigl({}^{(1)}\!\RootLiftMin{\sigma}(\ell)\bigr)
&\equiv
{}^{(2)}\!\RootLiftMin{\sigma}(\ell),
\label{eq:eqlift}
\\
\RootMapSub{\sigma}\bigl({}^{(1)}\!\RootSubMin{\sigma}(m)\bigr)
&\equiv
{}^{(2)}\!\RootSubMin{\sigma}(m),
\label{eq:eqsub}
\\
\RootMapShell{\sigma}
&\equiv
\bigl({}^{(2)}\!\RootProj{\sigma}\big|_{{}^{(2)}\!\RootShellSet{\sigma}}\bigr)^{-1}
\circ
{}^{(1)}\!\RootProj{\sigma}\big|_{{}^{(1)}\!\RootShellSet{\sigma}},
\label{eq:eqshell}
\\
\RootMapSupport{\sigma}
&\equiv
\bigl({}^{(2)}\!\RootProj{\sigma}\big|_{{}^{(2)}\!\RootSupportShell{\sigma}}\bigr)^{-1}
\circ
{}^{(1)}\!\RootProj{\sigma}\big|_{{}^{(1)}\!\RootSupportShell{\sigma}},
\label{eq:eqsupport}
\\
\RootMapZero{\sigma}
&\equiv
\bigl({}^{(2)}\!\RootProj{\sigma}\big|_{{}^{(2)}\!\RootZeroCore{\sigma}}\bigr)^{-1}
\circ
{}^{(1)}\!\RootProj{\sigma}\big|_{{}^{(1)}\!\RootZeroCore{\sigma}}.
\label{eq:eqzero}
\end{align}
The two cores are \emph{root-faithfully equivalent} if, under the maps \eqref{eq:eqfdn}--\eqref{eq:eqzero}, every displayed structure in Definition \ref{def:core} is transported identically.
\end{definition}

\begin{remark}
Definition \ref{def:equiv} is intentionally narrower than full off-shell identity of the ambient root body. The present benchmark-facing burden is not to classify every possible root thickening outside the fiber-minimizing exact-shell core. It is to remove the benchmark-relevant root freedom still visible on the active one-metric line.
\end{remark}

\section{Necessity of the Root Core}

\begin{proposition}[The fiber-minimizer hierarchy is necessary]
\label{prop:minimizers}
Every root-faithful presentation determines the same foundation-, ground-, origin-, precursor-, lift-, and substrate-fiber minimizer images inside the root layer. Hence the maps \eqref{eq:eqfdn}--\eqref{eq:eqsub} are well defined.
\end{proposition}

\begin{proof}
By Definition \ref{def:faithful}, the root presentation must recover exactly the fixed foundation package of Definition \ref{def:foundation}. In particular, it must recover the same foundation body and the same lower bodies together with the same unique minimizers on the corresponding recorded fibers. Since those lower fibers are fixed continuation data, the image of each root minimizer is keyed uniquely by the lower point it lifts. The comparison maps \eqref{eq:eqfdn}--\eqref{eq:eqsub} therefore depend only on the shared lower datum and are canonical.
\end{proof}

\begin{proposition}[Exact-shell transport data are necessary]
\label{prop:shell}
Every root-faithful presentation determines the same exact-shell transport data at benchmark scope. More precisely, the projection \(\RootProj{\sigma}\) restricts to diffeomorphisms from the root shell, support shell, and zero core onto the corresponding fixed foundation loci, and the root readouts, shell chart, support recorder, support-shell chart, zero chart, and zero/hidden split are thereby forced up to the canonical comparison maps \eqref{eq:eqshell}--\eqref{eq:eqzero}.
\end{proposition}

\begin{proof}
If a root presentation is faithful, then by definition its projected exact-shell transport recovers precisely the shell-side data already fixed in Definition \ref{def:foundation}. The root shell must therefore project diffeomorphically to the recorded foundation shell, otherwise the fixed foundation shell would not be recovered as exact-shell image. The same argument applies one step further to the support shell and then to the zero core. Once those three projected loci are fixed, the root observer-data and shell-response readouts, the shell chart to \(\ResSpace{\sigma}\), the support recorder, the support-shell chart to \(\EqnSpace{\sigma}\), the zero chart to \(\EqnZero{\sigma}\), and the zero/hidden split in the same equation variables cannot vary independently: changing any of them would change the projected foundation shell transport or the fixed equation-side coordinates already recorded on the active line. Hence all these constituents are forced up to transport by \eqref{eq:eqshell}--\eqref{eq:eqzero}.
\end{proof}

\begin{proposition}[Hidden charge, zero-response realization, solvability, and coercivity are necessary]
\label{prop:hidden}
Every root-faithful presentation determines, up to the canonical comparison maps, the same hidden charge data, the same exact root zero-response realization, the same fixed-data solvability statement, and the same coercive control on data-invisible directions.
\end{proposition}

\begin{proof}
The fixed foundation package already records one hidden-seed space \(\SeedHidden{\sigma}\), one hidden charge \(\FoundCharge{\sigma}\) with positive gap, one exact zero-response realization \(\FoundEmbed{\sigma}\), and one solvability/coercive package on the foundation zero core. A root-faithful presentation must project to that same foundation package. Because the support-shell split and the zero/hidden coordinates are already fixed by Proposition \ref{prop:shell}, the hidden sector cannot be changed without changing the induced foundation recorder. Likewise, the exact root zero-response realization must project to the same fixed visible, current, and equilibrium shell data, so its image on the root zero core is forced by the same foundation compatibility requirement. The solvability and coercive statements are properties of this same projected zero-response core; if they differed at root scope, the induced foundation statements would differ as well. Therefore those constituents are necessary up to the canonical comparison maps.
\end{proof}

\begin{proposition}[The benchmark package remains fixed]
\label{prop:benchmark}
Every root-faithful presentation preserves the same benchmark one-metric package \eqref{eq:fixedoutputs}. In particular, the operative metric remains exactly \(\CompletedMetric\), the ordinary matter route is unchanged, there is no second operative metric, and the low-energy matter law is not enlarged.
\end{proposition}

\begin{proof}
This is part of Definition \ref{def:faithful}. The present note removes only root-level freedom internal to the derivational line. It does not alter the fixed primitive outputs.
\end{proof}

\section{Main Theorem}

\begin{theorem}[Foundation-generating substrate root dynamics necessity / uniqueness]
\label{thm:main}
Fix the primitive continuation data of Definition \ref{def:fixed} and the recorded foundation package of Definition \ref{def:foundation}. Then, for each sector \(\sigma\in\{\Car,\Cur,\Hom\}\), the fiber-minimizing exact-shell root core \(\RootCore{\sigma}\) is necessary and unique at benchmark scope up to root-faithful equivalence. More precisely:
\begin{enumerate}
    \item every root-faithful presentation must determine the same fiber-minimizer hierarchy and the same exact-shell transport data described in Definition \ref{def:core};
    \item for any two root-faithful presentations, the canonical comparison maps \eqref{eq:eqfdn}--\eqref{eq:eqzero} are well defined and yield a unique root-faithful equivalence;
    \item under that equivalence, the root shell, support shell, zero core, and root readouts coincide, as do the shell/support/equation charts, the hidden charge, the exact zero-response realization, the fixed-data solvability statement, and the coercive control;
    \item consequently no additional benchmark-relevant root-level freedom remains on the active one-metric line beyond root-faithful relabeling.
\end{enumerate}
\end{theorem}

\begin{proof}
Item (1) is exactly the combined content of Propositions \ref{prop:minimizers}, \ref{prop:shell}, and \ref{prop:hidden}. Those propositions show that no root-faithful presentation may vary the fiber-minimizing exact-shell root core while still inducing the same fixed foundation package.

For item (2), Proposition \ref{prop:minimizers} gives the canonical comparison maps on the six minimizer images, while Proposition \ref{prop:shell} gives the canonical comparison maps on the shell, support shell, and zero core. Because each of those loci projects to the same fixed lower datum, the maps \eqref{eq:eqfdn}--\eqref{eq:eqzero} are unique.

For item (3), Proposition \ref{prop:shell} shows that the root readouts, shell chart, support recorder, support-shell chart, zero chart, and split are transported identically under the canonical comparison maps. Proposition \ref{prop:hidden} then shows that the hidden charge, exact zero-response realization, solvability statement, and coercive control are likewise transported identically. Hence the entire root core agrees under the unique root-faithful equivalence.

Item (4) is the project-level consequence of the previous items. Any remaining off-shell freedom outside the fiber-minimizing exact-shell root core is not benchmark relevant on the current line, because it does not change the fixed foundation package or the benchmark outputs. Therefore the current honest root burden has been removed in the only sense that matters for the active one-metric derivational program.
\end{proof}

\begin{corollary}[No remaining benchmark-relevant root choice]
\label{cor:nofreedom}
At current scope there is no second benchmark-relevant root presentation of the recorded foundation package other than a root-faithful relabeling of \(\RootCore{\sigma}\) in each sector.
\end{corollary}

\begin{proof}
If a second benchmark-relevant root presentation existed, it would either change some constituent proved necessary in Theorem \ref{thm:main}, contradicting item (1), or else differ only by the canonical identifications of item (2), in which case it would be merely a root-faithful relabeling.
\end{proof}

\begin{remark}
Corollary \ref{cor:nofreedom} is the precise benchmark-facing sense in which the present manuscript sharpens the recorded root derivation into a new frontier theorem. The current root layer is no longer merely available because a deeper theorem reproduces the fixed foundation package. The benchmark-relevant root core is now forced at current scope.
\end{remark}

\section{Routing Consequence}

\begin{corollary}[What must be done next]
\label{cor:next}
After Theorem \ref{thm:main}, the next honest primary manuscript below the current frontier must do at least one of the following:
\begin{enumerate}
    \item derive or justify the foundation-generating substrate root dynamics from still more primitive substrate structure;
    \item identify a genuinely smaller theorem-level core strictly inside \(\RootCore{\sigma}\) and prove a sharper collapse to it;
    \item do either of the previous two while preserving the same benchmark one-metric package, the same ordinary matter route, and the same claim boundary.
\end{enumerate}
It is no longer enough merely to restate the current root layer or to insert another same-output relay above it.
\end{corollary}

\begin{proof}
Theorem \ref{thm:main} removes the benchmark-relevant freedom presently visible at the root level. The next honest burden therefore lies either below that level or in a strictly smaller internal compression of it. Another package-preserving restatement of the same root core would not change project status.
\end{proof}

\section{Conclusion}

The positive result of this note is narrower than a completed first-principles derivation and stronger than another same-output deeper relay. The recorded root derivation theorem had already shown that the current foundation package is the canonical projected exact-shell output of a deeper root class. The present theorem then removes the remaining benchmark-relevant freedom inside that root layer itself. At current scope, the fiber-minimizing exact-shell root core is necessary and unique up to canonical root-faithful equivalence.

That conclusion is exactly the sharper theorem-level gain now needed below the former foundation frontier. The root-to-foundation projection, the six root fiber-minimizer images, the root shell and support transport, the zero/hidden split in the same equation variables, the hidden charge, the exact root zero-response realization, and the solvability/coercive package are no longer merely one possible deeper presentation of the same downstream line. Once the presentation is required to induce the fixed foundation package, those constituents are forced.

The scope remains disciplined. The benchmark package is unchanged: the one operative metric is still exactly \(\CompletedMetric\), the ordinary matter route is unchanged, there is no second operative metric, and the low-energy matter law is not enlarged. This remains a theorem inside the active derivational program of \TheoryName{}, not a basis for stronger promotion language. The next honest burden is deeper again: derive or justify the root dynamics from still more primitive substrate structure, or else prove a genuinely smaller collapse inside the current root core. Until then, adoption and derivation should continue to be kept sharply separated.

\end{document}
